equalities 2.4.2 are polynomials in @@M0@@, with rational coefficients, and to prove their equality it therefore suffices to show that they take the same value at all integer values of @@M1@@. But the relations obtained by substituting integer values for @@M2@@ are those expressing that the pre-@@M3@@-ring @@M4@@ (2.5.a)) is a @@M5@@-ring; they are therefore satisfied.

\sgaStableTarget{sga6.e05.def.2-7}\textbf{Definition 2.7.} \uline{By a} binomial ring \uline{we mean a ring satisfying conditions} 2.6.i) \uline{and} ii), \uline{endowed with the} @@M6@@-\uline{ring structure defined by} 2.6.1.

\sgaStableTarget{sga6.e05.ex.2-8}\textbf{Example 2.8.} Let @@M7@@ be a site, @@M8@@ a binomial ring, and @@M9@@ the sheaf associated to the constant presheaf with value @@M10@@ on @@M11@@; then @@M12@@ is a binomial ring.

Indeed, conditions 2.6 are preserved under passage to an inductive limit.

In particular, @@M13@@ is a binomial ring.

\sgaStableTarget{sga6.e05.para.2-9}2.9. Let @@M14@@ be a binomial ring, and @@M15@@ a @@M16@@-algebra. For every series @@M17@@, and every @@M18@@, we can define a series @@M19@@ in the usual way: we write @@M20@@, and put:
\sgaStableTarget{sga6.e05.eq.2-9-1}@@M21@@
the last equality being writable only modulo torsion. These formulas are obtained by substitution from the analogous formulas for @@M22@@; the following relations, the classical properties of the exponential, therefore remain valid for @@M23@@, @@M24@@:
\sgaStableTarget{sga6.e05.eq.2-9-2}@@M25@@

% Native English continuation macros retained at their source boundary.
\providecommand{\Ghat}{\widehat G}
\providecommand{\Gg}{\widehat G_g}
\providecommand{\Gcirc}{\widehat G_{\circ}}
\providecommand{\Gnabla}{\widehat G_{\nabla}}
\providecommand{\Grhat}{\widehat{\operatorname{Gr}}}
\providecommand{\Log}{\operatorname{Log}}
\providecommand{\Fil}{\operatorname{Fil}}
\providecommand{\eps}{\varepsilon}
\providecommand{\id}{\operatorname{id}}
\providecommand{\Sym}{\operatorname{Sym}}
\providecommand{\Mon}{\operatorname{Mon}}

% Exposé V continuation macros
\providecommand{\Ghat}{\widehat G}
\providecommand{\Gg}{\widehat G_g}
\providecommand{\Gcirc}{\widehat G_{\circ}}
\providecommand{\Gnabla}{\widehat G_{\nabla}}
\providecommand{\Grhat}{\widehat{\operatorname{Gr}}}
\providecommand{\Log}{\operatorname{Log}}
\providecommand{\Fil}{\operatorname{Fil}}
\providecommand{\eps}{\varepsilon}
\renewcommand{\id}{\operatorname{id}}
\providecommand{\Sym}{\operatorname{Sym}}
\providecommand{\Mon}{\operatorname{Mon}}

